%%%%%%%%%%%%%%%%%%%%%%%%%%%%%%%%%%%%%%%%%%%%%%%%%%%%%%%%%%%%%%%%%%%%%%%%%%%%%%%
%  SÍLABO DE ECONOMETRÍA — Plan de estudio integrado
%
%  Texto base:
%    DM  — Davidson y MacKinnon, Econometric Theory and Methods
%          (OUP, 2004; edición de los autores, 2021)
%  Textos de apoyo:
%    Wo  — Wooldridge, Introductory Econometrics (5.ª ed., 2012)
%          [trad.: Introducción a la econometría, 4.ª ed., 2009]
%    Hn  — Hansen, Econometrics (manuscrito, revisión de marzo de 2021)
%    Hy  — Hayashi, Econometrics (Princeton, 2000)
%    Gr  — Greene, Econometric Analysis (7.ª ed., 2012)
%  Consulta (sin horas): Davidson y MacKinnon (1993), Hansen (2022),
%    White (2001), Cameron y Trivedi (2005), Wooldridge (2002),
%    Angrist y Pischke (2009), Huber, Hastie, Tibshirani y Friedman,
%    Kleiber y Zeileis, Campbell, Lo y MacKinlay.
%  Programas PUCP integrados: 2ECO01 (2026-1) y ECO387 (2025-0).
%  Programas PUCP como material opcional: ECO387 (2022-2), ECO386,
%    ECO224, FIN214, FIN222, EST225.
%
%  Compilar con:  pdflatex Econometria.tex   (dos veces, por el índice)
%%%%%%%%%%%%%%%%%%%%%%%%%%%%%%%%%%%%%%%%%%%%%%%%%%%%%%%%%%%%%%%%%%%%%%%%%%%%%%%

\documentclass[11pt,a4paper]{article}

\usepackage[T1]{fontenc}
\usepackage[utf8]{inputenc}
\usepackage[spanish,es-noquoting,es-tabla]{babel}
\usepackage{lmodern}
\usepackage[scaled=0.92]{helvet}        % títulos en sans serif, al estilo McGraw-Hill
\usepackage[margin=2.4cm,headheight=15pt]{geometry}
\usepackage[table]{xcolor}
\usepackage{enumitem}
\usepackage{titlesec}
\usepackage{titletoc}
\usepackage{tabularx}
\usepackage{booktabs}
\usepackage{longtable}
\usepackage{amsmath}
\usepackage{amssymb}
\usepackage{fancyhdr}
\usepackage{tikz}
\usepackage[skins]{tcolorbox}
\usepackage{needspace}

%------------------------------------------------------------------ colores ---
% Gama de verdes
\definecolor{tonoOscuro}{RGB}{4,82,43}              % oscuro
\definecolor{tonoProfundo}{RGB}{6,121,64}           % profundo
\definecolor{tonoPrincipal}{RGB}{8,160,84}          % principal
\definecolor{tonoA}{RGB}{25,179,133}                % variante cálida
\definecolor{tonoB}{RGB}{15,179,69}                 % variante fría
\definecolor{tonoC}{RGB}{21,193,50}                 % variante viva
\definecolor{tonoClaro}{RGB}{95,221,158}            % claro
\definecolor{tonoAcero}{RGB}{63,131,97}             % acero
\definecolor{tonoPizarra}{RGB}{86,134,110}          % pizarra
\definecolor{tonoGrafito}{RGB}{39,73,56}            % grafito (texto secundario)
\definecolor{tonoHielo}{RGB}{207,247,227}           % hielo
\definecolor{tonoPalido}{RGB}{233,251,242}          % pálido
\definecolor{tonoBruma}{RGB}{246,254,250}           % bruma

\colorlet{azulosc}{tonoOscuro}
\colorlet{grisclaro}{tonoPalido}
% Un tono distinto de la gama para cada libro
\definecolor{colorDM}{RGB}{15,138,32}
\definecolor{colorWo}{RGB}{20,184,72}
\definecolor{colorHn}{RGB}{18,161,89}
\definecolor{colorHy}{RGB}{15,138,99}
\definecolor{colorGr}{RGB}{20,184,162}
\colorlet{colorPUCP}{tonoPizarra}

\usepackage[colorlinks=true,linkcolor=tonoOscuro,urlcolor=tonoPrincipal]{hyperref}

%------------------------------------------------------------------ estilos ---
\titleformat{\section}[block]
  {\sffamily\LARGE\bfseries\color{tonoOscuro}\raggedright}{\thesection.}{0.6em}
  {\textcolor{tonoPrincipal}{\rule[-0.1ex]{0.55em}{0.85em}}\hspace{0.45em}}
  [\vspace{-0.4ex}{\color{tonoPrincipal}\titlerule[1.6pt]}\vspace{1.2pt}{\color{tonoClaro}\titlerule[0.6pt]}]
\titlespacing*{\section}{0pt}{2.6ex plus 1ex minus .2ex}{1.6ex plus .2ex}

\titleformat{\subsection}
  {\sffamily\large\bfseries\color{tonoPrincipal}}{\thesubsection}{0.6em}
  {\textcolor{tonoClaro}{\rule[0.1ex]{0.45em}{0.45em}}\hspace{0.45em}}
\titlespacing*{\subsection}{0pt}{2.2ex plus .8ex minus .2ex}{1ex plus .2ex}

% Índice
\titlecontents{section}[0pt]
  {\addvspace{3pt}\sffamily\bfseries\color{tonoOscuro}}
  {}{\textcolor{tonoB}{\rule[0.15ex]{0.4em}{0.4em}}\hspace{0.6em}}
  {{\color{tonoClaro}\titlerule*[0.7em]{.}}\textcolor{tonoPrincipal}{\contentspage}}

% Cabecera y pie
\pagestyle{fancy}
\fancyhf{}
\lhead{\small\sffamily\bfseries\color{tonoOscuro}Sílabo de Econometría}
\rhead{\footnotesize\sffamily\color{tonoPrincipal}Davidson--MacKinnon $\cdot$ Wooldridge $\cdot$ Hansen $\cdot$ Hayashi $\cdot$ Greene}
\cfoot{{\setlength{\fboxsep}{3pt}\colorbox{tonoPrincipal}{\makebox[1.8em]{\small\sffamily\bfseries\color{white}\thepage}}}}
\renewcommand{\headrule}{%
  {\color{tonoPrincipal}\hrule height 1.1pt width\headwidth}\vskip 1.2pt
  {\color{tonoClaro}\hrule height 0.4pt width\headwidth}\vskip -2.7pt}

% Listas
\setlist[itemize,1]{label=\textcolor{tonoPrincipal}{\scriptsize$\blacksquare$}}
\setlist[enumerate,1]{label=\textcolor{tonoPrincipal}{\sffamily\bfseries\arabic*.}}

% Tablas: cabecera oscura y filas alternas en tono pálido
\arrayrulecolor{tonoPrincipal}
\setlength{\aboverulesep}{0pt}
\setlength{\belowrulesep}{0pt}
\renewcommand{\arraystretch}{1.3}
\newcommand{\cab}[1]{\textcolor{white}{\sffamily\bfseries #1}}
\newcommand{\filasalternas}{\rowcolors{2}{white}{tonoPalido}}

% Encabezados de párrafo en color
\newcommand{\destacado}[1]{{\sffamily\bfseries\color{tonoPrincipal}#1}}

% Franja decorativa con la gama
\newcommand{\franjas}[1][3pt]{%
  \noindent\begin{tikzpicture}
    \foreach \c [count=\i from 0] in {tonoOscuro,tonoProfundo,tonoPrincipal,tonoA,tonoB,tonoC,tonoClaro,tonoHielo}
      \fill[\c] (\i*\linewidth/8,0) rectangle ++(\linewidth/8,#1);
  \end{tikzpicture}\par}

%---------------------------------------------------------------- comandos ----
% Etiqueta de libro: recuadro de color con la sigla
\newcommand{\etiqueta}[2]{%
  {\setlength{\fboxsep}{1.4pt}\colorbox{#1}{\sffamily\bfseries\color{white}\strut#2}}}
% Referencias abreviadas a cada libro
\newcommand{\DM}[1]{\etiqueta{colorDM}{DM}\nobreak\hspace{0.25em}\textcolor{colorDM}{#1}}
\newcommand{\WO}[1]{\etiqueta{colorWo}{Wo}\nobreak\hspace{0.25em}\textcolor{colorWo}{#1}}
\newcommand{\HN}[1]{\etiqueta{colorHn}{Hn}\nobreak\hspace{0.25em}\textcolor{colorHn}{#1}}
\newcommand{\HY}[1]{\etiqueta{colorHy}{Hy}\nobreak\hspace{0.25em}\textcolor{colorHy}{#1}}
\newcommand{\GR}[1]{\etiqueta{colorGr}{Gr}\nobreak\hspace{0.25em}\textcolor{colorGr}{#1}}

% Un tema: número en recuadro de color + título
\newcommand{\tema}[1]{\temaaux#1\relax}
\def\temaaux#1 #2\relax{%
  \par\addvspace{1.1em}\noindent
  {\sffamily\bfseries\large\raggedright
   \hangindent=2.9em\hangafter=1
   \makebox[2.9em][l]{{\setlength{\fboxsep}{2.5pt}\colorbox{tonoPrincipal}{\color{white}#1}}}%
   {\color{tonoOscuro}#2}\par}%
  \nobreak\vspace{0.25em}\noindent{\color{tonoClaro}\rule{\linewidth}{0.6pt}}\par\nobreak\vspace{0.35em}}

% Línea de lecturas con marcador y etiqueta de color
\newcommand{\linealectura}[3]{%
  \par\nobreak\noindent\hspace{1.2em}%
  \begin{minipage}{0.92\textwidth}\small
    \textcolor{#1}{\rule[0.1ex]{0.45em}{0.45em}}\hspace{0.45em}%
    {\sffamily\bfseries\color{#1}#2}\enspace#3%
  \end{minipage}\par\vspace{0.3em}}
\newcommand{\lecturas}[1]{\linealectura{tonoPrincipal}{Lecturas:}{#1}}
\newcommand{\apoyo}[1]{\linealectura{tonoA}{Profundización:}{#1}}
\newcommand{\juegos}[1]{\linealectura{tonoB}{Teoría de juegos:}{#1}}
\newcommand{\supl}[1]{\linealectura{tonoC}{Suplemento:}{#1}}
\newtcolorbox{cajanota}{enhanced, sharp corners, boxrule=0pt, frame hidden,
  colback=tonoPalido, borderline west={2.5pt}{0pt}{tonoClaro},
  left=8pt, right=6pt, top=4pt, bottom=4pt, before skip=0pt, after skip=0pt,
  fontupper=\small\color{tonoGrafito}}
\newcommand{\nota}[1]{%
  \par\nobreak\noindent\hspace{1.2em}%
  \begin{minipage}{0.92\textwidth}\begin{cajanota}#1\end{cajanota}\end{minipage}\par\vspace{0.3em}}
\newcommand{\consultalib}[1]{\linealectura{tonoPizarra}{Consulta:}{{\color{tonoGrafito}#1}}}
\newcommand{\consulta}[1]{\linealectura{colorPUCP}{Consulta (PUCP):}{{\color{tonoGrafito}#1}}}

% Encabezado de bloque: banda oscura, franja de la gama y recuadro de presentación
\newif\ifrecienparte
\newcommand{\bloque}[2]{%
  \ifrecienparte\global\recienpartefalse\else\par\needspace{0.4\textheight}\fi
  \phantomsection\addcontentsline{toc}{section}{#1}%
  \begin{tcolorbox}[enhanced, sharp corners, boxrule=0pt, frame hidden,
      colback=tonoOscuro, borderline west={7pt}{0pt}{tonoClaro},
      left=14pt, right=10pt, top=10pt, bottom=10pt, before skip=0pt, after skip=0pt]
    {\raggedright\sffamily\Large\bfseries\color{white}#1\par}
  \end{tcolorbox}%
  \nointerlineskip\franjas[4pt]
  \vspace{0.8em}
  \begin{tcolorbox}[enhanced, sharp corners, boxrule=0pt, frame hidden,
      colback=tonoPalido, borderline west={3pt}{0pt}{tonoPrincipal},
      borderline south={0.6pt}{0pt}{tonoClaro},
      left=10pt, right=8pt, top=6pt, bottom=6pt, before skip=0pt, after skip=0pt,
      fontupper=\small\color{tonoGrafito}]
    #2
  \end{tcolorbox}%
  \par\medskip}

%================================================================== documento ==
\begin{document}

\begin{titlepage}
\begin{tikzpicture}[remember picture,overlay]
  % banda superior
  \fill[tonoOscuro] (current page.north west) rectangle ([yshift=-5.4cm]current page.north east);
  \fill[tonoClaro] (current page.north west) rectangle ([xshift=1.1cm,yshift=-5.4cm]current page.north west);
  % franja de la gama bajo la banda
  \foreach \c [count=\i from 0] in {tonoOscuro,tonoProfundo,tonoPrincipal,tonoA,tonoB,tonoC,tonoClaro,tonoHielo}
    \fill[\c] ([xshift=\i*\paperwidth/8,yshift=-5.4cm]current page.north west)
      rectangle ++(\paperwidth/8,-0.35cm);
  % franja inferior
  \foreach \c [count=\i from 0] in {tonoHielo,tonoClaro,tonoC,tonoB,tonoA,tonoPrincipal,tonoProfundo,tonoOscuro}
    \fill[\c] ([xshift=\i*\paperwidth/8]current page.south west)
      rectangle ++(\paperwidth/8,0.6cm);
\end{tikzpicture}%
\centering
\vspace*{0.5cm}
{\sffamily\Huge\bfseries\color{white} Sílabo de Econometría\par}
\vspace{0.4cm}
{\sffamily\Large\color{tonoHielo} Plan de estudio integrado\par}
\vspace{1.9cm}
{\sffamily\large\bfseries\color{tonoPrincipal}\textcolor{tonoClaro}{\rule[0.5ex]{1.5cm}{0.6pt}}\quad Texto base\quad\textcolor{tonoClaro}{\rule[0.5ex]{1.5cm}{0.6pt}}\par}
\vspace{0.3cm}
{\large\etiqueta{colorDM}{DM} \quad {\color{tonoOscuro}Russell Davidson y James G. MacKinnon}\par}
{\normalsize \emph{\color{colorDM}Econometric Theory and Methods}\par}
{\small\color{tonoPizarra} Oxford University Press, 2004; edición de los autores, 2021 --- 15 capítulos\par}
\vspace{0.7cm}
{\sffamily\large\bfseries\color{tonoPrincipal}\textcolor{tonoClaro}{\rule[0.5ex]{1.5cm}{0.6pt}}\quad Textos de apoyo\quad\textcolor{tonoClaro}{\rule[0.5ex]{1.5cm}{0.6pt}}\par}
\vspace{0.3cm}
{\large\etiqueta{colorWo}{Wo} \quad {\color{tonoOscuro}Jeffrey M. Wooldridge}\par}
{\normalsize \emph{\color{colorWo}Introductory Econometrics: A Modern Approach}\par}
{\small\color{tonoPizarra} Cengage, 5.ª edición, 2012 --- 19 capítulos y 7 apéndices\par}
{\small\color{tonoPizarra} (traducción: \emph{Introducción a la econometría. Un enfoque moderno}, 4.ª ed., 2009)\par}
\vspace{0.3cm}
{\large\etiqueta{colorHn}{Hn} \quad {\color{tonoOscuro}Bruce E. Hansen}\par}
{\normalsize \emph{\color{colorHn}Econometrics}\par}
{\small\color{tonoPizarra} Manuscrito, revisión de marzo de 2021 (Princeton University Press, 2022) --- 29 capítulos\par}
\vspace{0.3cm}
{\large\etiqueta{colorHy}{Hy} \quad {\color{tonoOscuro}Fumio Hayashi}\par}
{\normalsize \emph{\color{colorHy}Econometrics}\par}
{\small\color{tonoPizarra} Princeton University Press, 2000 --- 10 capítulos\par}
\vspace{0.3cm}
{\large\etiqueta{colorGr}{Gr} \quad {\color{tonoOscuro}William H. Greene}\par}
{\normalsize \emph{\color{colorGr}Econometric Analysis}\par}
{\small\color{tonoPizarra} Pearson, 7.ª edición, 2012 --- 21 capítulos y 6 apéndices\par}
\vspace{0.7cm}
{\normalsize\color{colorPUCP} Con los temas de dos programas de la PUCP:\par}
{\normalsize\color{colorPUCP} Econometría Intermedia (2ECO01, maestría) y Tópicos de Econometría (ECO387)\par}
\vfill
{\setlength{\fboxsep}{6pt}\colorbox{tonoPrincipal}{\normalsize\sffamily\bfseries\color{white} Trece bloques temáticos y un pretramo --- 8 meses}\par}
\vspace{0.2cm}
\end{titlepage}

\tableofcontents
\newpage

%------------------------------------------------------------------------------
\section*{Cómo usar este sílabo}
\addcontentsline{toc}{section}{Cómo usar este sílabo}

El plan está organizado en \textbf{trece bloques temáticos} y un pretramo, y
sigue el orden de \textcolor{colorDM}{\textbf{Davidson y MacKinnon}}: primero el
modelo de regresión lineal y su geometría, luego la teoría asintótica y la
simulación, después los estimadores para modelos más generales (mínimos
cuadrados no lineales, generalizados, variables instrumentales, método
generalizado de momentos y máxima verosimilitud) y, al final, los sistemas de
ecuaciones, las variables dependientes limitadas y la especificación. Conviene
cerrar cada bloque antes de abrir el siguiente.

\medskip
\noindent\destacado{Criterio general de reparto} (en el orden en que se lee cada tema):
\begin{itemize}[leftmargin=1.4em,itemsep=2pt]
  \item \textcolor{colorWo}{\textbf{Wooldridge}} aporta la intuición y los
        ejemplos empíricos. Se lee \emph{primero}, cuando el tema lo tiene,
        para entender qué problema se resuelve. No cubre GMM, máxima
        verosimilitud general, simulación ni métodos no paramétricos.
  \item \textcolor{colorDM}{\textbf{Davidson y MacKinnon}} es el texto base y
        el que fija el orden. Se lee \emph{completo} en cada tema que lo cita:
        su enfoque geométrico y sus regresiones artificiales unifican todo el
        plan.
  \item \textcolor{colorHn}{\textbf{Hansen}} es la referencia moderna y la
        lectura principal donde Davidson y MacKinnon son breves o no llegan:
        esperanza condicional y proyección, teoría asintótica, bootstrap,
        errores agrupados, instrumentos débiles, regresión no paramétrica y
        aprendizaje automático. Cubre los temas del curso ECO387, que sigue el
        primer año del doctorado de Chicago.
  \item \textcolor{colorHy}{\textbf{Hayashi}} es la lectura principal de GMM
        (Bloque 9) y de los estimadores extremos (tema 6.4), y una alternativa
        rigurosa en MCO, asintótica y máxima verosimilitud.
  \item \textcolor{colorGr}{\textbf{Greene}} es la referencia enciclopédica:
        aplicaciones, variantes de cada modelo y apéndices de matemáticas y
        estadística. Se lee \emph{después}, o se consulta, según el tema.
\end{itemize}

\medskip
\noindent\destacado{Qué entra en las horas.} Las líneas \emph{Lecturas} y
\emph{Profundización} entran en el cómputo de horas. \emph{Profundización}
señala la lectura de mayor rigor (casi siempre Hansen o Hayashi): se aborda
cuando el tema ya se entendió con los textos anteriores. Las líneas
\emph{Consulta} y \emph{Consulta (PUCP)} remiten a libros de la carpeta y a
referencias de los programas que no entran en el cómputo de horas (Anexo E).

\medskip
\noindent\destacado{Material excluido.} Quedan fuera las series de tiempo
(\DM{13--14}, \HY{6}, \HY{9--10}, \HN{14--16}, \GR{21}, \WO{10--11},
\WO{18}) y los datos de panel (\DM{7.10}, \HY{5}, \HN{17--18}, \GR{11},
\WO{13--14}), que tienen sus propios materiales en la carpeta. La correlación
serial se estudia solo como violación de un supuesto del modelo de regresión
(tema 7.3). También quedan fuera los métodos bayesianos y los modelos de
duración. El detalle está en el Anexo B.

\medskip
\noindent\destacado{Los vacíos de los programas.} Ninguno de los dos programas de
la PUCP trata a fondo el método generalizado de momentos, el bootstrap, los
errores agrupados ni la estimación no paramétrica. El plan los incorpora con
los libros de la carpeta: el Bloque 5 (simulación y bootstrap), el tema 7.4
(errores agrupados), el Bloque 9 (GMM) y el Bloque 13 (especificación,
métodos no paramétricos y aprendizaje automático).

\medskip
\noindent\destacado{Sobre las matemáticas.} El pretramo repasa álgebra
matricial, probabilidad e inferencia estadística, que se usan desde el
Bloque 1. Para la teoría asintótica del Bloque 4 conviene tener a mano
\GR{Apéndice D} y el capítulo 4 de Davidson y MacKinnon (1993).

\medskip
\noindent\destacado{Sobre los ejercicios y el software.} La econometría se
aprende derivando resultados y estimando modelos con datos. Davidson y
MacKinnon incluyen ejercicios con datos al final de cada capítulo; Hansen,
ejercicios empíricos y secciones de programación; Wooldridge, ejercicios de
computadora con sus bases de datos. Conviene hacer cada semana al menos un
ejercicio teórico y uno con datos, en R, Python o Stata.

\medskip
\noindent\destacado{Notación.} \DM{n} remite al capítulo $n$ de Davidson y
MacKinnon, y \WO{n}, \HN{n}, \HY{n} y \GR{n}, al capítulo de los demás
textos. Si solo entra parte de un capítulo se indican secciones: \DM{8.3} es
la sección 3 del capítulo 8. La numeración de Wooldridge es la de la 5.ª
edición en inglés; la traducción de la carpeta (4.ª edición) conserva los
mismos capítulos, aunque alguna sección puede variar. La de Hansen es la del
manuscrito de 2021 que está en la carpeta.

%------------------------------------------------------------------------------
\newpage
\section*{Cronograma por semanas}
\addcontentsline{toc}{section}{Cronograma por semanas}

Calendario del plan: 340 horas distribuidas en \textbf{34 semanas de
contenido}, que equivalen a \textbf{unos 8 meses de calendario}, a un ritmo de
diez horas semanales.

\medskip
\noindent Todas las semanas son de materia nueva y no hay semanas de colchón.
La columna de horas es la que gobierna el reparto. Si hay semanas perdidas por
feriados, viajes o imprevistos, el ajuste correcto es alargar el calendario,
no comprimir los bloques.

\medskip
{\filasalternas
\noindent\begin{tabularx}{\textwidth}{lXcc}
\toprule
\rowcolor{tonoOscuro}\cab{Semanas} & \cab{Contenido} & \cab{Sem.} & \cab{Horas} \\
\midrule
1--2 & Pretramo --- Álgebra matricial, probabilidad e inferencia & 2 & 20 \\
3--4 & Bloque 1 --- Esperanza condicional, proyección y el modelo de regresión & 2 & 20 \\
5--6 & Bloque 2 --- Mínimos cuadrados: álgebra y geometría & 2 & 20 \\
7--8 & Bloque 3 --- MCO en muestras finitas e inferencia exacta & 2 & 20 \\
\rowcolor{tonoHielo}\textbf{9--11} & \textbf{Bloque 4 --- Teoría asintótica e inferencia en muestras grandes} & \textbf{3} & \textbf{30} \\
12--13 & Bloque 5 --- Simulación, Monte Carlo y bootstrap & 2 & 20 \\
14--15 & Bloque 6 --- Regresión no lineal y estimadores extremos & 2 & 20 \\
16--17 & Bloque 7 --- Mínimos cuadrados generalizados y errores no esféricos & 2 & 20 \\
\rowcolor{tonoHielo}\textbf{18--20} & \textbf{Bloque 8 --- Variables instrumentales} & \textbf{3} & \textbf{30} \\
\rowcolor{tonoHielo}\textbf{21--23} & \textbf{Bloque 9 --- Método generalizado de momentos} & \textbf{3} & \textbf{30} \\
\rowcolor{tonoHielo}\textbf{24--26} & \textbf{Bloque 10 --- Máxima verosimilitud} & \textbf{3} & \textbf{30} \\
27--28 & Bloque 11 --- Sistemas de ecuaciones y ecuaciones simultáneas & 2 & 20 \\
\rowcolor{tonoHielo}\textbf{29--31} & \textbf{Bloque 12 --- Variables dependientes discretas y limitadas} & \textbf{3} & \textbf{30} \\
\rowcolor{tonoHielo}\textbf{32--34} & \textbf{Bloque 13 --- Especificación, selección de modelos y métodos no paramétricos} & \textbf{3} & \textbf{30} \\
\midrule
\rowcolor{tonoHielo}\multicolumn{2}{l}{\textbf{Total}} & \textbf{34} & \textbf{340} \\
\bottomrule
\end{tabularx}
}

\medskip
\noindent\destacado{Ritmo supuesto.} Diez horas por semana, todas de contenido
nuevo: por ejemplo, cuatro sesiones de dos horas entre semana para la teoría y
una sesión de dos horas el sábado para ejercicios con datos.

\medskip
\noindent\destacado{Los tramos largos.} El Bloque 4 es el más importante: la
teoría asintótica sostiene todo lo que viene después. Los Bloques 8 a 10
(semanas 18--26) forman un recorrido continuo por los tres grandes principios
de estimación ---variables instrumentales, momentos y verosimilitud--- y son
el punto donde el avance suele atascarse.

\medskip
\noindent\destacado{Hitos de control.} Conviene verificar el avance al terminar las
semanas 11, 20, 26 y 31. Si en alguno de esos cortes se acumula más de una
semana de retraso, se alarga el calendario y no se comprime el bloque
siguiente.

\subsection*{Detalle semana a semana}

{\small\filasalternas
\renewcommand{\arraystretch}{1.1}\noindent\begin{tabularx}{\textwidth}{clXc}
\toprule
\rowcolor{tonoOscuro}\cab{Sem.} & \cab{Bloque} & \cab{Temas} & \cab{Horas} \\
\midrule
1  & Pretramo & P.1 Álgebra matricial; P.2 Probabilidad y distribuciones & 10 \\
2  & Pretramo & P.3 Inferencia estadística & 10 \\
\midrule
3  & 1 & 1.1 Econometría, datos y causalidad; 1.2 Esperanza condicional & 10 \\
4  & 1 & 1.3 Mejor predictor lineal; 1.4 Interpretaciones del modelo de regresión & 10 \\
\midrule
5  & 2 & 2.1 Método de momentos y MCO; 2.2 Geometría de MCO & 10 \\
6  & 2 & 2.3 Teorema FWL; 2.4 Observaciones influyentes; 2.5 Forma funcional y ajuste & 10 \\
\midrule
7  & 3 & 3.1 Insesgadez, varianza y Gauss--Markov; 3.2 Mala especificación & 10 \\
8  & 3 & 3.3 Inferencia exacta en el modelo normal; 3.4 Estabilidad de parámetros & 10 \\
\midrule
9  & 4 & 4.1 Convergencia, ley de los grandes números y teorema central del límite & 10 \\
10 & 4 & 4.2 Asintótica de MCO; 4.3 Covarianzas robustas y método delta & 10 \\
11 & 4 & 4.4 Pruebas e intervalos asintóticos; 4.5 Estimación restringida; 4.6 Potencia & 10 \\
\midrule
12 & 5 & 5.1 Experimentos de Monte Carlo; 5.2 El bootstrap & 10 \\
13 & 5 & 5.3 Intervalos y pruebas bootstrap; 5.4 Jackknife, bootstrap salvaje y por conglomerados & 10 \\
\midrule
14 & 6 & 6.1 Mínimos cuadrados no lineales; 6.2 Cómputo y optimización numérica & 10 \\
15 & 6 & 6.3 Regresión de Gauss--Newton; 6.4 Estimadores extremos & 10 \\
\midrule
16 & 7 & 7.1 MCG y MCGF; 7.2 Heteroscedasticidad & 10 \\
17 & 7 & 7.3 Correlación serial en el modelo lineal; 7.4 Errores agrupados & 10 \\
\midrule
18 & 8 & 8.1 Endogeneidad; 8.2 Variables instrumentales y MC2E & 10 \\
19 & 8 & 8.3 Propiedades de los estimadores VI; 8.4 Contrastes de endogeneidad y sobreidentificación & 10 \\
20 & 8 & 8.5 Instrumentos débiles y LIML; 8.6 LATE y experimentos naturales & 10 \\
\midrule
21 & 9 & 9.1 GMM lineal; 9.2 GMM eficiente y matrices HAC & 10 \\
22 & 9 & 9.3 Contrastes en GMM; 9.4 GMM con varias ecuaciones & 10 \\
23 & 9 & 9.5 GMM no lineal, distancia mínima y momentos simulados & 10 \\
\midrule
24 & 10 & 10.1 Conceptos e identificación; 10.2 Propiedades asintóticas & 10 \\
25 & 10 & 10.3 Matriz de información; 10.4 Las tres pruebas clásicas & 10 \\
26 & 10 & 10.5 Cuasi-verosimilitud y estimación en dos etapas; 10.6 Transformaciones & 10 \\
\midrule
27 & 11 & 11.1 Regresión multivariada y SUR; 11.2 Identificación en ecuaciones simultáneas & 10 \\
28 & 11 & 11.3 MC3E, FIML y LIML & 10 \\
\midrule
29 & 12 & 12.1 Respuesta binaria & 10 \\
30 & 12 & 12.2 Respuestas múltiples; 12.3 Datos de conteo & 10 \\
31 & 12 & 12.4 Truncamiento y censura; 12.5 Selección muestral & 10 \\
\midrule
32 & 13 & 13.1 Pruebas de especificación; 13.2 Hipótesis no anidadas y criterios de información & 10 \\
33 & 13 & 13.3 Regresión no paramétrica; 13.4 Regresión por series; 13.5 Regresión cuantílica & 10 \\
34 & 13 & 13.6 Contracción, \emph{lasso} y aprendizaje automático & 10 \\
\bottomrule
\end{tabularx}
}

%------------------------------------------------------------------------------
\newpage
\bloque{Pretramo --- Álgebra matricial, probabilidad e inferencia --- 2 semanas}
{No es un bloque de contenido econométrico sino la puesta a punto de las
herramientas que el plan da por supuestas desde el Bloque 1. El curso ECO387
supone el nivel de Casella y Berger; aquí se alcanza con los apéndices de
Greene y Wooldridge y, si hace falta, con Hansen (2022).}

\tema{P.1 Álgebra matricial}
\lecturas{\WO{Apéndice D} \quad \DM{1.4} \quad \GR{Apéndice A}}
\consultalib{Davidson y MacKinnon (1993), apéndice A}
\nota{Rango, inversa, matrices particionadas, formas cuadráticas, matrices
definidas positivas, matrices idempotentes y valores propios. Las matrices de
proyección del Bloque 2 se apoyan en todo esto.}

\tema{P.2 Probabilidad, momentos y distribuciones}
\lecturas{\WO{Apéndice B} \quad \DM{1.2} \quad \GR{Apéndice B}}
\consultalib{Hansen (2022), caps.~1--5 \quad Davidson y MacKinnon (1993), apéndice B}
\nota{Distribuciones conjuntas, marginales y condicionales; esperanza y
varianza condicionales; normal multivariada y sus distribuciones derivadas
($\chi^2$, $t$, $F$).}

\tema{P.3 Inferencia estadística: estimación, pruebas e intervalos}
\lecturas{\WO{Apéndice C} \quad \GR{Apéndice C}}
\consultalib{Hansen (2022), caps.~6, 10, 11, 13 y 14}
\consulta{Casella y Berger (2021) \quad Lehmann y Romano (2005)}
\nota{Estimadores y sus propiedades (sesgo, eficiencia, consistencia),
método de momentos y máxima verosimilitud en su versión elemental, pruebas de
hipótesis e intervalos de confianza. Los libros de Casella y Berger y de
Lehmann y Romano, que recomienda ECO387, no están en la carpeta.}

%------------------------------------------------------------------------------
\bloque{Bloque 1 --- Esperanza condicional, proyección y el modelo de regresión --- 2 semanas}
{El punto de partida de ECO387: antes de estimar, definir qué parámetro se
quiere conocer. La regresión aparece primero como esperanza condicional o
como proyección lineal (uso descriptivo) y solo después, bajo supuestos
adicionales, como relación causal o estructural.}

\tema{1.1 Econometría, datos y causalidad}
\lecturas{\WO{1} \quad \DM{1.1} \quad \GR{1}}
\apoyo{\HN{1}}
\consultalib{Angrist y Pischke (2009), caps.~1--2 \quad Cameron y Trivedi (2005), cap.~2}
\nota{Estructura de los datos económicos, el ideal experimental y la noción de
\emph{ceteris paribus}.}

\tema{1.2 La esperanza condicional y su error}
\lecturas{\HN{2.1--2.17} \quad \DM{1.3}}
\nota{Ley de esperanzas iteradas, error de la esperanza condicional, mejor
predictor, varianza condicional, homoscedasticidad y heteroscedasticidad,
derivada de la regresión. Hansen es aquí la lectura principal: es el tema 2
de ECO387.}

\tema{1.3 Mejor predictor lineal y proyección}
\lecturas{\HN{2.18--2.29}}
\nota{La proyección lineal existe siempre, aunque la esperanza condicional no
sea lineal. Sesgo de variable omitida en la proyección, regresión hacia la
media y límites de la proyección lineal.}

\tema{1.4 Interpretaciones del modelo de regresión: descriptiva, causal y estructural}
\lecturas{\WO{2.1} \quad \WO{3.1} \quad \GR{2}}
\apoyo{\HN{2.30}}
\consultalib{Angrist y Pischke (2009), cap.~3}
\nota{Los supuestos del modelo lineal clásico (\GR{2.3}) y lo que se gana con
cada uno. Es el tema 3 de ECO387 y el tema 1 de la Unidad 1 de 2ECO01.}

%------------------------------------------------------------------------------
\bloque{Bloque 2 --- Mínimos cuadrados: álgebra y geometría --- 2 semanas}
{MCO como cálculo antes que como estimador. La geometría de Davidson y
MacKinnon (proyecciones ortogonales, teorema de Frisch--Waugh--Lovell) es la
herramienta que se reutiliza en todo el plan: en MCG, en VI, en las
regresiones artificiales y en las pruebas de especificación.}

\tema{2.1 Método de momentos y estimador MCO}
\lecturas{\WO{2.2--2.3} \quad \WO{3.2} \quad \WO{Apéndice E} \quad \DM{1.5} \quad \GR{3.1--3.2}}
\apoyo{\HN{3.1--3.9}}

\tema{2.2 La geometría de MCO}
\lecturas{\DM{2.1--2.3}}
\apoyo{\HN{3.10--3.16} \quad \HY{1.2}}
\consultalib{Davidson y MacKinnon (1993), cap.~1}
\nota{Espacios vectoriales, matrices de proyección $P_X$ y aniquiladora $M_X$,
residuos como proyección ortogonal.}

\tema{2.3 El teorema de Frisch--Waugh--Lovell y sus aplicaciones}
\lecturas{\DM{2.4--2.5} \quad \GR{3.3--3.4}}
\apoyo{\HN{3.17--3.18}}
\nota{Regresión particionada, variables centradas, tendencias y variables
binarias estacionales.}

\tema{2.4 Observaciones influyentes y \emph{leverage}}
\lecturas{\DM{2.6}}
\apoyo{\HN{3.19--3.21} \quad \HN{3.24}}
\nota{Valores de \emph{leverage}, regresión dejando una observación fuera y
colinealidad numérica.}

\tema{2.5 Bondad de ajuste, forma funcional y variables binarias}
\lecturas{\WO{6} \quad \WO{7.1--7.4} \quad \DM{3.8} \quad \GR{3.5} \quad \GR{6.2--6.3}}
\nota{Escalamiento de datos, logaritmos, términos cuadráticos e
interacciones, $R^2$ y $\bar{R}^2$. Es la parte de «variables cualitativas»
de la Unidad 1 de 2ECO01.}

%------------------------------------------------------------------------------
\bloque{Bloque 3 --- MCO en muestras finitas e inferencia exacta --- 2 semanas}
{Las propiedades de MCO cuando el modelo es correcto y lo que ocurre cuando no
lo es. Cubre los temas 1 a 3 de la Unidad 1 de 2ECO01.}

\tema{3.1 Insesgadez, matriz de covarianzas y teorema de Gauss--Markov}
\lecturas{\WO{3.3--3.5} \quad \DM{3.1--3.5} \quad \GR{4.1--4.3}}
\apoyo{\HY{1.1} \quad \HY{1.3} \quad \HN{4.1--4.9} \quad \HN{4.12--4.15}}
\nota{Hansen incluye la versión moderna del teorema de Gauss--Markov
(\HN{4.9}), que elimina la restricción a estimadores lineales.}

\tema{3.2 Residuos, mala especificación, variables omitidas y multicolinealidad}
\lecturas{\DM{3.6--3.7} \quad \WO{9.2} \quad \GR{4.7}}
\apoyo{\HN{4.22}}
\nota{Sesgo de variable omitida e inclusión de variables irrelevantes;
multicolinealidad imperfecta y sus síntomas. Es el tema 3 de la Unidad 1 de
2ECO01.}

\tema{3.3 Inferencia exacta en el modelo lineal normal}
\lecturas{\WO{4} \quad \DM{4.1--4.4} \quad \GR{5.1--5.5}}
\apoyo{\HY{1.4--1.5} \quad \HN{5}}
\nota{Distribuciones $\chi^2$, $t$ y $F$; pruebas $t$ y $F$; restricciones
lineales; relación con máxima verosimilitud bajo normalidad.}

\tema{3.4 Estabilidad de parámetros y pruebas de especificación sencillas}
\lecturas{\WO{7.4} \quad \WO{9.1} \quad \GR{6.4} \quad \GR{5.9}}
\consulta{García (2022), \emph{Econometría 1}}
\nota{Prueba de Chow, estimación recursiva y CUSUM (pedidos por 2ECO01) y
prueba RESET de Ramsey.}

%------------------------------------------------------------------------------
\bloque{Bloque 4 --- Teoría asintótica e inferencia en muestras grandes --- 3 semanas}
{El bloque más importante del plan y el tema 1 de ECO387. Sin normalidad de
los errores, toda la inferencia descansa en resultados de muestras grandes:
consistencia, normalidad asintótica y estimación robusta de covarianzas.
Todo lo que sigue (VI, GMM, verosimilitud) reutiliza estas herramientas.}

\tema{4.1 Modos de convergencia, ley de los grandes números y teorema central del límite}
\lecturas{\HN{6} \quad \GR{Apéndice D}}
\apoyo{\HY{2.1}}
\consultalib{Davidson y MacKinnon (1993), cap.~4 \quad White (2001), caps.~2--5 \quad Hansen (2022), caps.~7--9}
\consulta{van der Vaart (2000)}
\nota{Convergencia en probabilidad y en distribución, teorema de la aplicación
continua, método delta, notación $O_p$ y $o_p$. Davidson y MacKinnon (1993)
dedica un capítulo completo a estos resultados, que el texto base resume.}

\tema{4.2 Consistencia y normalidad asintótica de MCO}
\lecturas{\WO{5} \quad \DM{3.3} \quad \GR{4.4}}
\apoyo{\HN{7.1--7.5} \quad \HY{2.3}}

\tema{4.3 Covarianzas robustas a heteroscedasticidad y método delta}
\lecturas{\WO{8.1--8.2} \quad \DM{5.5--5.6}}
\apoyo{\HN{7.6--7.11} \quad \HY{2.5--2.6}}
\consultalib{White (2001), cap.~6}
\nota{Estimadores HC0 a HC3, errores estándar de funciones no lineales de los
parámetros.}

\tema{4.4 Pruebas e intervalos asintóticos}
\lecturas{\DM{4.5} \quad \DM{5.1--5.2} \quad \DM{5.4} \quad \GR{5.6--5.7}}
\apoyo{\HN{7.12--7.18} \quad \HN{9.1--9.16} \quad \HY{2.4}}
\nota{Pruebas de Wald, $t$ asintóticas, regiones de confianza y restricciones
no lineales.}

\tema{4.5 Estimación restringida y distancia mínima}
\lecturas{\HN{8}}
\nota{MCO restringido, estimador de distancia mínima eficiente y restricciones
de desigualdad. Prepara el terreno para el GMM del Bloque 9.}

\tema{4.6 Potencia, valores $p$ y uso de las pruebas}
\lecturas{\DM{4.7}}
\apoyo{\HN{9.17--9.23}}
\nota{Errores de tipo I y II, potencia asintótica local, pruebas múltiples y
corrección de Bonferroni.}

%------------------------------------------------------------------------------
\bloque{Bloque 5 --- Simulación, Monte Carlo y bootstrap --- 2 semanas}
{\textbf{Bloque nuevo respecto de los programas de la PUCP.} 2ECO01 usa un
experimento de Monte Carlo para ilustrar la consistencia, pero ninguno de los
dos programas estudia el bootstrap. Es la herramienta estándar para hacer
inferencia cuando la teoría asintótica es poco fiable.}

\tema{5.1 Experimentos de Monte Carlo}
\lecturas{\DM{4.6} \quad \GR{15.1--15.2} \quad \GR{15.5}}
\apoyo{\HN{9.18}}
\consultalib{Davidson y MacKinnon (1993), caps.~21.1--21.7 \quad Cameron y Trivedi (2005), cap.~12}
\nota{Generación de números pseudoaleatorios, diseño de experimentos y
técnicas de reducción de varianza.}

\tema{5.2 El bootstrap: principio, varianzas y errores estándar}
\lecturas{\HN{10.1--10.15} \quad \GR{15.4}}
\consultalib{Cameron y Trivedi (2005), cap.~11}

\tema{5.3 Intervalos y pruebas bootstrap}
\lecturas{\DM{5.3}}
\apoyo{\HN{10.16--10.28}}
\consultalib{Davidson y MacKinnon (1993), cap.~21.8}
\nota{Intervalos percentil, BC, BCa y percentil-$t$; refinamientos
asintóticos; pruebas bootstrap y bootstrap paramétrico.}

\tema{5.4 Jackknife, bootstrap salvaje y bootstrap por conglomerados}
\lecturas{\HN{10.3--10.5} \quad \HN{10.29--10.30}}
\nota{El bootstrap salvaje es la versión válida con heteroscedasticidad; el
bootstrap por conglomerados enlaza con el tema 7.4.}

%------------------------------------------------------------------------------
\bloque{Bloque 6 --- Regresión no lineal y estimadores extremos --- 2 semanas}
{La extensión de MCO a modelos no lineales y el marco general de los
estimadores definidos como solución de un problema de optimización. Es el
puente teórico hacia GMM (Bloque 9) y máxima verosimilitud (Bloque 10).}

\tema{6.1 Estimadores de momentos no lineales y mínimos cuadrados no lineales}
\lecturas{\DM{6.1--6.3} \quad \GR{7.1--7.2}}
\apoyo{\HN{23.1--23.5}}
\consultalib{Davidson y MacKinnon (1993), caps.~2 y 5}

\tema{6.2 Cómputo y optimización numérica}
\lecturas{\DM{6.4} \quad \GR{Apéndice E}}
\apoyo{\HY{7.5} \quad \HN{23.9}}
\nota{Métodos de Newton y cuasi-Newton, criterios de convergencia y
problemas prácticos.}

\tema{6.3 La regresión de Gauss--Newton y las pruebas en modelos no lineales}
\lecturas{\DM{6.5--6.8}}
\consultalib{Davidson y MacKinnon (1993), cap.~6}
\nota{Las regresiones artificiales son el sello de Davidson y MacKinnon:
reaparecen en las pruebas de especificación del Bloque 13.}

\tema{6.4 Estimadores extremos: consistencia y normalidad asintótica}
\lecturas{\HY{7.1--7.4}}
\apoyo{\HN{22} \quad \GR{12}}
\consultalib{Wooldridge (2002), cap.~12}
\nota{Identificación, ley uniforme de los grandes números y normalidad
asintótica de un estimador M. Hayashi es aquí la lectura principal.}

%------------------------------------------------------------------------------
\bloque{Bloque 7 --- Mínimos cuadrados generalizados y errores no esféricos --- 2 semanas}
{Las violaciones de los supuestos sobre la matriz de covarianzas de los
errores: heteroscedasticidad, correlación serial y dependencia dentro de
grupos. Cubre los temas 1 y 2 de la Unidad 2 de 2ECO01.}

\tema{7.1 Mínimos cuadrados generalizados y factibles}
\lecturas{\DM{7.1--7.4} \quad \GR{9.1--9.3}}
\apoyo{\HY{1.6} \quad \HN{4.10--4.11}}

\tema{7.2 Heteroscedasticidad: pruebas y mínimos cuadrados ponderados}
\lecturas{\WO{8.3--8.5} \quad \DM{7.5} \quad \GR{9.4--9.7}}
\apoyo{\HY{2.7--2.8}}
\nota{Pruebas de Goldfeld--Quandt, Breusch--Pagan y White; MCP y el modelo de
probabilidad lineal.}

\tema{7.3 Correlación serial en el modelo de regresión}
\lecturas{\WO{12.1--12.5} \quad \DM{7.6--7.9} \quad \GR{20.5} \quad \GR{20.7--20.9}}
\apoyo{\HY{2.10}}
\nota{Solo como violación de un supuesto del modelo lineal: pruebas de
Durbin--Watson y Breusch--Godfrey, Cochrane--Orcutt y Prais--Winsten. Los
procesos ARMA y la teoría de series de tiempo quedan fuera del plan.}

\tema{7.4 Errores agrupados (\emph{clustered})}
\lecturas{\HN{4.23--4.25}}
\consultalib{Angrist y Pischke (2009), cap.~8 \quad Cameron y Trivedi (2005), cap.~24}
\consulta{Cameron y Miller (2015) \quad Abadie, Athey, Imbens y Wooldridge (2017)}
\nota{\textbf{Tema nuevo respecto de los programas.} Inferencia con muestras
agrupadas y a qué nivel agrupar. Es la práctica habitual con encuestas como la
ENAHO (Anexo D.2).}

%------------------------------------------------------------------------------
\bloque{Bloque 8 --- Variables instrumentales --- 3 semanas}
{El tema 3 de la Unidad 2 de 2ECO01 y el tema 5 de ECO387. Hansen amplía el
tratamiento con los problemas que dominan la práctica actual: instrumentos
débiles, muchos instrumentos y la interpretación LATE.}

\tema{8.1 Endogeneidad: variables omitidas, error de medida y simultaneidad}
\lecturas{\WO{15.1} \quad \WO{15.4} \quad \WO{9.4} \quad \DM{8.1--8.2} \quad \GR{8.1--8.2} \quad \GR{8.5}}
\apoyo{\HN{12.1--12.6} \quad \HY{3.1--3.2}}

\tema{8.2 El estimador de variables instrumentales y MC2E}
\lecturas{\WO{15.2--15.3} \quad \DM{8.3} \quad \GR{8.3}}
\apoyo{\HN{12.7--12.12}}
\nota{Forma reducida, condiciones de identificación (orden y rango),
estimador de Wald y mínimos cuadrados en dos etapas.}

\tema{8.3 Propiedades de los estimadores VI en muestras grandes y finitas}
\lecturas{\DM{8.4}}
\apoyo{\HN{12.15--12.22} \quad \HN{12.26}}
\consultalib{White (2001), cap.~1}
\nota{Consistencia y normalidad asintótica de MC2E, estimación de su
covarianza, sesgo en muestras finitas y regresores generados.}

\tema{8.4 Contrastes de endogeneidad y de restricciones de sobreidentificación}
\lecturas{\WO{15.5--15.6} \quad \DM{8.5--8.8} \quad \GR{8.4}}
\apoyo{\HN{12.28--12.33}}
\nota{Pruebas de Durbin--Wu--Hausman, función de control, prueba de Sargan y
pruebas bootstrap (\DM{8.8}).}

\tema{8.5 Instrumentos débiles, muchos instrumentos y LIML}
\lecturas{\HN{12.35--12.39} \quad \GR{8.7}}
\apoyo{\HN{12.13--12.14} \quad \HN{12.19}}
\nota{\textbf{Tema nuevo respecto de los programas.} Identificación débil,
prueba de Stock y Yogo, el estimador LIML y el JIVE.}

\tema{8.6 Efectos de tratamiento locales y experimentos naturales}
\lecturas{\HN{12.34} \quad \HN{12.40--12.41} \quad \GR{8.8}}
\consultalib{Angrist y Pischke (2009), cap.~4 \quad Huber, cap.~6}
\nota{LATE y las dos aplicaciones clásicas que replica Hansen: Acemoglu,
Johnson y Robinson (2001), que también es lectura de ECO386, y Angrist y
Krueger (1991).}

%------------------------------------------------------------------------------
\bloque{Bloque 9 --- Método generalizado de momentos --- 3 semanas}
{\textbf{Bloque nuevo respecto de los programas de la PUCP}: ninguno lo trata a
fondo. GMM unifica MCO, VI, MC2E, SUR y MC3E como casos particulares y es el
método de estimación de buena parte de la macroeconomía empírica. Hayashi es
la lectura principal.}

\tema{9.1 El método de momentos y GMM en modelos lineales}
\lecturas{\DM{9.1--9.2} \quad \HY{3.3--3.5} \quad \GR{13.1--13.2} \quad \GR{13.4}}
\apoyo{\HN{13.1--13.9}}

\tema{9.2 GMM eficiente, matriz de ponderación y covarianzas HAC}
\lecturas{\DM{9.3} \quad \HY{3.5} \quad \HY{3.8}}
\apoyo{\HN{13.10--13.13} \quad \HN{13.20}}
\nota{GMM en dos etapas, iterado y de actualización continua (CUE); GMM bajo
homoscedasticidad condicional como MC2E.}

\tema{9.3 Contrastes en GMM}
\lecturas{\DM{9.4} \quad \HY{3.6--3.7} \quad \GR{13.5}}
\apoyo{\HN{13.14--13.24}}
\nota{Prueba $J$ de Hansen, pruebas de distancia, pruebas de subconjuntos de
condiciones y de endogeneidad en el marco de GMM.}

\tema{9.4 GMM con varias ecuaciones}
\lecturas{\HY{4.1--4.6}}
\nota{FIVE, MC3E y SUR como casos especiales de GMM multiecuacional;
prepara el Bloque 11.}

\tema{9.5 GMM no lineal, distancia mínima y método de momentos simulados}
\lecturas{\DM{9.5--9.6} \quad \GR{13.3} \quad \GR{13.6}}
\apoyo{\HN{13.25--13.27}}
\consultalib{Wooldridge (2002), cap.~14 \quad Cameron y Trivedi (2005), cap.~6}
\nota{Modelos de momentos condicionales y la estimación de modelos con
expectativas racionales.}

%------------------------------------------------------------------------------
\bloque{Bloque 10 --- Máxima verosimilitud --- 3 semanas}
{El tercer gran principio de estimación y el tema 6 de ECO387. Se estudia como
caso particular de los estimadores extremos del tema 6.4, lo que permite
reutilizar los resultados de consistencia y normalidad asintótica.}

\tema{10.1 Conceptos básicos e identificación}
\lecturas{\DM{10.1--10.2} \quad \GR{14.1--14.3}}
\consultalib{Davidson y MacKinnon (1993), cap.~8 \quad Hansen (2022), cap.~10}

\tema{10.2 Propiedades asintóticas del estimador de máxima verosimilitud}
\lecturas{\DM{10.3} \quad \GR{14.4}}
\apoyo{\HY{7.2--7.3}}
\nota{Consistencia, normalidad asintótica y eficiencia: la cota de
Cramér--Rao.}

\tema{10.3 La matriz de información y la covarianza del estimador}
\lecturas{\DM{10.4}}
\nota{Igualdad de la matriz de información; estimadores por hessiana,
producto externo del gradiente (OPG) y sándwich.}

\tema{10.4 Las tres pruebas clásicas: razón de verosimilitud, Wald y multiplicador de Lagrange}
\lecturas{\DM{10.5--10.6} \quad \GR{14.6}}
\apoyo{\HY{7.4}}
\consultalib{Davidson y MacKinnon (1993), cap.~13}
\nota{Su equivalencia asintótica (tema 4 de la Unidad 1 de 2ECO01) y su
cálculo mediante regresiones artificiales.}

\tema{10.5 Cuasi-verosimilitud, estimación en dos etapas y relación con GMM}
\lecturas{\GR{14.5} \quad \GR{14.7--14.8}}
\consultalib{Wooldridge (2002), cap.~13}
\nota{Qué estima un estimador de máxima verosimilitud cuando el modelo está mal
especificado, y covarianzas robustas.}

\tema{10.6 Transformaciones de la variable dependiente}
\lecturas{\DM{10.8}}
\consultalib{Davidson y MacKinnon (1993), cap.~14}
\nota{El jacobiano en la verosimilitud y la transformación de Box--Cox.}

%------------------------------------------------------------------------------
\bloque{Bloque 11 --- Sistemas de ecuaciones y ecuaciones simultáneas --- 2 semanas}
{Los modelos multiecuacionales de la sumilla de 2ECO01 (recursivos, SUR y
ecuaciones simultáneas) y el tema 4 de su Unidad 2. Se ubica después de GMM
y de máxima verosimilitud porque los estimadores de información completa se
apoyan en ambos.}

\tema{11.1 Regresión multivariada y regresiones aparentemente no relacionadas (SUR)}
\lecturas{\DM{12.1--12.3} \quad \GR{10.1--10.5}}
\apoyo{\HN{11.1--11.10}}

\tema{11.2 Ecuaciones simultáneas: forma estructural, forma reducida e identificación}
\lecturas{\WO{16.1--16.4} \quad \DM{12.4} \quad \GR{10.6}}
\nota{Sesgo de simultaneidad, condiciones de orden y rango, modelos
recursivos.}

\tema{11.3 Estimación de información limitada y completa: MC3E, FIML y LIML}
\lecturas{\DM{12.5--12.6} \quad \DM{12.8} \quad \HY{8.5--8.6}}
\apoyo{\HY{4.5}}
\consultalib{Wooldridge (2002), caps.~7--9}

%------------------------------------------------------------------------------
\bloque{Bloque 12 --- Variables dependientes discretas y limitadas --- 3 semanas}
{La Unidad 3 de 2ECO01. Es la aplicación más importante de máxima
verosimilitud en microeconometría. Los modelos de duración quedan fuera.}

\tema{12.1 Modelos de respuesta binaria: MPL, logit y probit}
\lecturas{\WO{7.5} \quad \WO{17.1} \quad \DM{11.1--11.3} \quad \GR{17.1--17.3}}
\apoyo{\HN{25.1--25.10} \quad \HY{8.1}}
\consultalib{Wooldridge (2002), cap.~15 \quad Cameron y Trivedi (2005), cap.~14}
\nota{Variable latente, efectos marginales, bondad de ajuste y pruebas
mediante la regresión binaria artificial (\DM{11.3}).}

\tema{12.2 Respuestas múltiples: probit bivariado, logit multinomial, ordenado y anidado}
\lecturas{\DM{11.4} \quad \GR{17.5} \quad \GR{18.2--18.3}}
\apoyo{\HN{26}}
\consultalib{Cameron y Trivedi (2005), cap.~15}

\tema{12.3 Modelos para datos de conteo}
\lecturas{\WO{17.3} \quad \DM{11.5} \quad \GR{18.4}}
\consultalib{Wooldridge (2002), cap.~19 \quad Cameron y Trivedi (2005), cap.~20}

\tema{12.4 Truncamiento, censura y el modelo tobit}
\lecturas{\WO{17.2} \quad \WO{17.4} \quad \DM{11.6} \quad \GR{19.1--19.3}}
\apoyo{\HN{27.1--27.8} \quad \HY{8.2--8.3}}
\consultalib{Wooldridge (2002), cap.~16 \quad Cameron y Trivedi (2005), cap.~16}
\consulta{Amemiya (1985), cap.~10}

\tema{12.5 Selección muestral: el modelo de Heckman}
\lecturas{\WO{17.5} \quad \DM{11.7} \quad \GR{19.5}}
\apoyo{\HN{27.9--27.11}}
\consultalib{Wooldridge (2002), cap.~17}
\nota{Estimación en dos etapas con la inversa de la razón de Mills y por
máxima verosimilitud con información completa.}

%------------------------------------------------------------------------------
\bloque{Bloque 13 --- Especificación, selección de modelos y métodos no paramétricos --- 3 semanas}
{Cierra el plan con dos vacíos de los programas: cómo evaluar un modelo y qué
hacer cuando no se quiere imponer una forma funcional. El último tema abre la
puerta al aprendizaje automático, que el curso ECO386 conecta con la
inferencia causal (Anexo D.1).}

\tema{13.1 Pruebas de especificación con regresiones artificiales}
\lecturas{\DM{15.1--15.2} \quad \DM{15.7} \quad \GR{5.9}}
\consultalib{Davidson y MacKinnon (1993), caps.~11--12}
\nota{Pruebas de momentos condicionales, de la matriz de información y de
normalidad de los errores.}

\tema{13.2 Hipótesis no anidadas y criterios de información}
\lecturas{\WO{6.3} \quad \DM{15.3--15.4} \quad \GR{5.8} \quad \GR{5.10}}
\apoyo{\HN{28.1--28.14}}
\nota{Pruebas $J$ y $P$, criterios de Akaike y de Schwarz, validación
cruzada y estrategia de lo general a lo particular.}

\tema{13.3 Regresión no paramétrica}
\lecturas{\DM{15.5} \quad \HN{19}}
\apoyo{\GR{7.5} \quad \GR{12.4}}
\nota{Estimadores de núcleo y polinomios locales, elección del ancho de
banda, maldición de la dimensionalidad.}

\tema{13.4 Regresión por series y modelos parcialmente lineales}
\lecturas{\HN{20} \quad \HN{19.24} \quad \GR{7.4}}
\apoyo{\GR{12.3}}
\consultalib{Cameron y Trivedi (2005), cap.~9}

\tema{13.5 Regresión cuantílica}
\lecturas{\WO{9.6} \quad \GR{7.3}}
\apoyo{\HN{24}}
\consultalib{Angrist y Pischke (2009), cap.~7}
\nota{Desviaciones absolutas mínimas y cuantiles condicionales.}

\tema{13.6 Selección de modelos, contracción, \emph{lasso} y aprendizaje automático}
\lecturas{\HN{28.15--28.31} \quad \HN{29}}
\consultalib{Hastie, Tibshirani y Friedman (2009) \quad Huber, cap.~5}
\nota{Inferencia después de seleccionar un modelo, estimadores de
James--Stein, promedios de modelos, \emph{ridge}, \emph{lasso}, árboles y
bosques aleatorios, \emph{lasso} doble y aprendizaje automático
doble/insesgado (DML). Es la base para el Anexo D.1 y para el trabajo empírico
sobre los efectos de la IA.}

%------------------------------------------------------------------------------
\newpage
\section*{Anexo A. De los capítulos a los bloques}
\addcontentsline{toc}{section}{Anexo A. De los capítulos a los bloques}

Capítulos de cada libro que entran en el plan y temas donde se leen. Las
secciones excluidas se detallan en el Anexo B.

\newcommand{\mapa}[3]{%
\subsection*{#1}
{\small\filasalternas
\noindent\begin{tabularx}{\textwidth}{>{\raggedright\arraybackslash}p{1.1cm}X>{\raggedright\arraybackslash}p{4.0cm}}
\toprule
\rowcolor{tonoOscuro}\cab{Cap.} & \cab{Título} & \cab{Temas} \\
\midrule
#3
\bottomrule
\end{tabularx}
}}

\mapa{A.1 Davidson y MacKinnon, \emph{Econometric Theory and Methods} (texto base)}{}{%
1 & Regression Models & P.1, P.2, 1.1, 1.2, 2.1 \\
2 & The Geometry of Linear Regression & 2.2, 2.3, 2.4 \\
3 & The Statistical Properties of Ordinary Least Squares & 2.5, 3.1, 3.2, 4.2 \\
4 & Hypothesis Testing in Linear Regression Models & 3.3, 4.4, 4.6, 5.1 \\
5 & Confidence Intervals & 4.3, 4.4, 5.3 \\
6 & Nonlinear Regression & 6.1, 6.2, 6.3 \\
7 & Generalized Least Squares and Related Topics (sin 7.10) & 7.1, 7.2, 7.3 \\
8 & Instrumental Variables Estimation & 8.1--8.4 \\
9 & The Generalized Method of Moments & 9.1, 9.2, 9.3, 9.5 \\
10 & The Method of Maximum Likelihood (sin 10.7) & 10.1--10.4, 10.6 \\
11 & Discrete and Limited Dependent Variables (sin 11.8) & 12.1--12.5 \\
12 & Multivariate Models & 11.1, 11.2, 11.3 \\
15 & Testing the Specification of Econometric Models & 13.1, 13.2, 13.3 \\
}

\mapa{A.2 Wooldridge, \emph{Introductory Econometrics}}{}{%
1 & The Nature of Econometrics and Economic Data & 1.1 \\
2 & The Simple Regression Model & 1.4, 2.1 \\
3 & Multiple Regression Analysis: Estimation & 1.4, 2.1, 3.1 \\
4 & Multiple Regression Analysis: Inference & 3.3 \\
5 & Multiple Regression Analysis: OLS Asymptotics & 4.2 \\
6 & Multiple Regression Analysis: Further Issues & 2.5, 13.2 \\
7 & Multiple Regression Analysis with Qualitative Information & 2.5, 3.4, 12.1 \\
8 & Heteroskedasticity & 4.3, 7.2 \\
9 & More on Specification and Data Issues & 3.2, 3.4, 8.1, 13.5 \\
12 & Serial Correlation and Heteroskedasticity (12.1--12.5) & 7.3 \\
15 & Instrumental Variables Estimation and Two Stage Least Squares (15.1--15.6) & 8.1, 8.2, 8.4 \\
16 & Simultaneous Equations Models (16.1--16.4) & 11.2 \\
17 & Limited Dependent Variable Models and Sample Selection Corrections & 12.1, 12.3, 12.4, 12.5 \\
Ap. B--E & Probabilidad, estadística, álgebra matricial y el modelo en forma matricial & P.1--P.3, 2.1 \\
}

\mapa{A.3 Hansen, \emph{Econometrics}}{}{%
1 & Introduction & 1.1 \\
2 & Conditional Expectation and Projection & 1.2, 1.3, 1.4 \\
3 & The Algebra of Least Squares & 2.1--2.4 \\
4 & Least Squares Regression & 3.1, 3.2, 7.1, 7.4 \\
5 & Normal Regression & 3.3 \\
6 & A Review of Large Sample Asymptotics & 4.1 \\
7 & Asymptotic Theory for Least Squares & 4.2, 4.3, 4.4 \\
8 & Restricted Estimation & 4.5 \\
9 & Hypothesis Testing & 4.4, 4.6, 5.1 \\
10 & Resampling Methods & 5.2, 5.3, 5.4 \\
11 & Multivariate Regression (11.1--11.10) & 11.1 \\
12 & Instrumental Variables & 8.1--8.6 \\
13 & Generalized Method of Moments & 9.1, 9.2, 9.3, 9.5 \\
19 & Nonparametric Regression & 13.3, 13.4 \\
20 & Series Regression & 13.4 \\
22 & M-Estimators & 6.4 \\
23 & Nonlinear Least Squares (23.1--23.5 y 23.9) & 6.1, 6.2 \\
24 & Quantile Regression & 13.5 \\
25 & Binary Choice (25.1--25.10) & 12.1 \\
26 & Multiple Choice & 12.2 \\
27 & Censoring and Selection (27.1--27.11) & 12.4, 12.5 \\
28 & Model Selection, Stein Shrinkage, and Model Averaging & 13.2, 13.6 \\
29 & Machine Learning & 13.6 \\
}

\mapa{A.4 Hayashi, \emph{Econometrics}}{}{%
1 & Finite-Sample Properties of OLS & 2.2, 3.1, 3.3, 7.1 \\
2 & Large-Sample Theory (2.1, 2.3--2.8, 2.10) & 4.1--4.4, 7.2, 7.3 \\
3 & Single-Equation GMM & 8.1, 9.1, 9.2, 9.3 \\
4 & Multiple-Equation GMM (4.1--4.6) & 9.4, 11.3 \\
7 & Extremum Estimators & 6.2, 6.4, 10.2, 10.4 \\
8 & Examples of Maximum Likelihood (8.1--8.3, 8.5--8.6) & 11.3, 12.1, 12.4 \\
}

\mapa{A.5 Greene, \emph{Econometric Analysis}}{}{%
1 & Econometrics & 1.1 \\
2 & The Linear Regression Model & 1.4 \\
3 & Least Squares & 2.1, 2.3, 2.5 \\
4 & The Least Squares Estimator & 3.1, 3.2, 4.2 \\
5 & Hypothesis Tests and Model Selection & 3.3, 3.4, 4.4, 13.1, 13.2 \\
6 & Functional Form and Structural Change & 2.5, 3.4 \\
7 & Nonlinear, Semiparametric, and Nonparametric Regression Models & 6.1, 13.3, 13.4, 13.5 \\
8 & Endogeneity and Instrumental Variable Estimation & 8.1--8.6 \\
9 & The Generalized Regression Model and Heteroscedasticity & 7.1, 7.2 \\
10 & Systems of Equations & 11.1, 11.2 \\
12 & Estimation Frameworks in Econometrics & 6.4, 13.3, 13.4 \\
13 & Minimum Distance Estimation and the GMM Estimator & 9.1, 9.2, 9.3, 9.5 \\
14 & Maximum Likelihood Estimation (14.1--14.8) & 10.1, 10.2, 10.4, 10.5 \\
15 & Simulation-Based Estimation and Inference (15.1--15.2, 15.4--15.5) & 5.1, 5.2 \\
17 & Discrete Choice (17.1--17.3, 17.5) & 12.1, 12.2 \\
18 & Discrete Choices and Event Counts & 12.2, 12.3 \\
19 & Limited Dependent Variables (19.1--19.3, 19.5) & 12.4, 12.5 \\
20 & Serial Correlation (20.5, 20.7--20.9) & 7.3 \\
Ap. A--E & Álgebra, probabilidad, inferencia, asintótica y optimización & P.1--P.3, 4.1, 6.2 \\
}

%------------------------------------------------------------------------------
\section*{Anexo B. Capítulos y secciones que no entran en el plan}
\addcontentsline{toc}{section}{Anexo B. Capítulos y secciones que no entran en el plan}

{\small\filasalternas
\noindent\begin{tabularx}{\textwidth}{>{\raggedright\arraybackslash}p{4.3cm}X}
\toprule
\rowcolor{tonoOscuro}\cab{Materia} & \cab{Capítulos y secciones} \\
\midrule
Series de tiempo (carpeta propia) & \DM{13--14}, \DM{10.7}; \WO{10--11}, \WO{12.6}, \WO{15.7}, \WO{16.5}, \WO{18}; \HN{14--16}; \HY{2.11--2.12}, \HY{6}, \HY{8.7}, \HY{9--10}; \GR{20} (salvo 20.5 y 20.7--20.9), \GR{21} \\
Datos de panel (carpeta propia) & \DM{7.10}; \WO{13--14}, \WO{15.8}, \WO{16.6}; \HN{17--18}, \HN{23.6}, \HN{25.13}, \HN{27.12}; \HY{5}; \GR{11}, \GR{17.4} \\
Métodos bayesianos & \GR{16} \\
Modelos de duración & \DM{11.8}; \GR{19.4} \\
Diseños de evaluación de impacto & \HN{21} (regresión discontinua); \GR{19.6}. Se recomiendan en el Anexo D.1 \\
Otras secciones & \HN{11.11--11.17} (rango reducido y modelos de factores); \HN{23.7--23.8} (modelos de umbral); \HN{25.11--25.12}; \HY{2.9}, \HY{4.7}, \HY{8.4}; \GR{14.9--14.10}, \GR{15.3}, \GR{15.6} en adelante \\
Proyecto empírico & \WO{19}: no entra en las horas; se usa en el Anexo D.6 \\
\bottomrule
\end{tabularx}
}

\medskip
\noindent Davidson y MacKinnon (1993), \emph{Estimation and Inference in
Econometrics}, queda como libro de consulta: cubre los mismos temas que el
texto base con más detalle técnico, y sus capítulos se citan en las líneas
«Consulta».

%------------------------------------------------------------------------------
\section*{Anexo C. Programas de la PUCP utilizados}
\addcontentsline{toc}{section}{Anexo C. Programas de la PUCP utilizados}

El plan se construye sobre dos programas de la carpeta. Los programas de
microeconometría (1ECO25, 1ECO77, ECO315) y de macroeconometría (1ECO24,
1ECO78) no se usaron.

\subsection*{C.1 Econometría Intermedia (2ECO01, Maestría en Economía, 2026-1)}
{\small\filasalternas
\noindent\begin{tabularx}{\textwidth}{>{\raggedright\arraybackslash}p{5.6cm}X}
\toprule
\rowcolor{tonoOscuro}\cab{Unidad y tema del programa} & \cab{Temas del plan} \\
\midrule
U1. Modelo clásico de regresión lineal y MCO & 1.4, 2.1, 2.2, 3.1 \\
U1. Evaluación de MCO, variables cualitativas y cambio estructural & 2.5, 3.3, 3.4 \\
U1. Multicolinealidad y variables omitidas & 3.2, 3.4 \\
U1. Muestras grandes, Monte Carlo y máxima verosimilitud & 4.2, 5.1, Bloque 10 \\
U2. Perturbaciones no esféricas y heteroscedasticidad & 7.1, 7.2 \\
U2. Autocorrelación & 7.3 \\
U2. Endogeneidad, VI y MC2E & 8.1--8.4 \\
U2. Ecuaciones simultáneas & Bloque 11 \\
U3. Respuesta binaria y probit bivariado & 12.1, 12.2 \\
U3. Extensiones multinomiales & 12.2 \\
U3. Truncamiento, censura y tobit & 12.4 \\
U3. Selección muestral & 12.5 \\
U4. Series de tiempo y predicción & Fuera del plan (carpeta Time Series) \\
\bottomrule
\end{tabularx}
}

\subsection*{C.2 Tópicos de Econometría (ECO387, 2025-0, R.~Quineche)}
Basado en el primer año del doctorado de la Universidad de Chicago.

\medskip
{\small\filasalternas
\noindent\begin{tabularx}{\textwidth}{>{\raggedright\arraybackslash}p{5.6cm}X}
\toprule
\rowcolor{tonoOscuro}\cab{Tema del programa} & \cab{Temas del plan} \\
\midrule
1. Teoría de muestras grandes & Bloque 4 \\
2. Esperanza condicional & 1.2, 1.3 \\
3. Interpretaciones del modelo de regresión lineal & 1.4 \\
4. Mínimos cuadrados ordinarios & Bloques 2 a 4 \\
5. Variables instrumentales & Bloque 8 \\
6. Estimadores de máxima verosimilitud & Bloques 6 y 10 \\
\bottomrule
\end{tabularx}
}

\subsection*{C.3 Vacíos cubiertos con los libros de la carpeta}
{\small\filasalternas
\noindent\begin{tabularx}{\textwidth}{>{\raggedright\arraybackslash}p{5.6cm}X}
\toprule
\rowcolor{tonoOscuro}\cab{Vacío} & \cab{Dónde se cubre} \\
\midrule
Método generalizado de momentos & Bloque 9 (Hayashi, Davidson y MacKinnon, Hansen) \\
Simulación y bootstrap & Bloque 5 (Hansen, Davidson y MacKinnon, Greene) \\
Errores agrupados & Tema 7.4 (Hansen) \\
Instrumentos débiles y LATE & Temas 8.5 y 8.6 (Hansen) \\
Mínimos cuadrados no lineales y estimadores extremos & Bloque 6 (Davidson y MacKinnon, Hayashi) \\
Especificación y métodos no paramétricos & Bloque 13 (Davidson y MacKinnon, Hansen, Greene) \\
\bottomrule
\end{tabularx}
}

%------------------------------------------------------------------------------
\section*{Anexo D. Material opcional}
\addcontentsline{toc}{section}{Anexo D. Material opcional}

Todo lo de este anexo es opcional: no entra en el cómputo de horas ni modifica
el cronograma. Sigue los demás programas de econometría de la carpeta y sirve
para ampliar el plan cuando el ritmo lo permita.

\subsection*{D.1 Inferencia causal y aprendizaje automático}
Sigue el programa ECO386 (Tópicos de Economía Matemática y Computacional,
2024-2, A.~Quispe), que une los métodos de identificación causal con el
aprendizaje automático. Momento sugerido: después del Bloque 13. Es el módulo
más cercano al trabajo empírico sobre los efectos de la IA en la economía.

\medskip
{\small\filasalternas
\noindent\begin{tabularx}{\textwidth}{>{\raggedright\arraybackslash}p{4.4cm}X>{\raggedright\arraybackslash}p{2.6cm}}
\toprule
\rowcolor{tonoOscuro}\cab{Tema} & \cab{Lecturas} & \cab{Relación con el plan} \\
\midrule
Experimentos aleatorios y efectos de tratamiento & Huber, caps.~2--3; Imbens y Rubin (2015) & 1.1, 1.4 \\
Selección en observables, \emph{matching} y ponderación & Huber, cap.~4; Angrist y Pischke (2009), cap.~3; \GR{19.6} & 1.4, 12.5 \\
Regresión en alta dimensión y \emph{lasso} doble & \HN{29.19--29.21}; Chernozhukov, Hansen y Spindler (2015) & 13.6 \\
DAG y modelos de ecuaciones estructurales & Heckman y Pinto (2015) & 11.2 \\
Aprendizaje automático doble/insesgado (DML) & \HN{29.22}; Chernozhukov et al. (2018a); Huber, cap.~5 & 13.6 \\
Efectos heterogéneos en experimentos & Chernozhukov, Demirer, Duflo y Fernández-Val (2018b) & 13.6 \\
DML con variables instrumentales y LATE & Huber, cap.~6 & 8.6 \\
Regresión discontinua y diferencias en diferencias & \HN{21}; Angrist y Pischke (2009), caps.~5--6; Huber, caps.~7 y 9 & 8.6 \\
Aprendizaje automático para economistas & Athey e Imbens (2019); Hastie, Tibshirani y Friedman (2009) & 13.6 \\
\emph{Embeddings} y precios hedónicos & Bajari et al. (2021) & 13.6 \\
\bottomrule
\end{tabularx}
}

\subsection*{D.2 Muestras complejas y errores agrupados}
Sigue el programa ECO387 (2022-2, A.~Cozzubo y J.~M.~Del Pozo). Las encuestas
del INEI (ENAHO, ENDES) tienen diseños con estratos, conglomerados y factores
de expansión; ignorarlos sesga los estimadores y subestima los errores
estándar. Momento sugerido: después del Bloque 7.

\begin{itemize}[leftmargin=1.4em,itemsep=1pt]
  \item Muestreo aleatorio simple, estratificado, por conglomerados y diseños
        complejos: Lumley (2010); Groves et al. (2011).
  \item Pesos muestrales en la regresión: Solon, Haider y Wooldridge (2015);
        Cameron y Trivedi (2005), cap.~24.
  \item Errores agrupados y bootstrap salvaje: Cameron y Miller (2015); Abadie,
        Athey, Imbens y Wooldridge (2017); \HN{4.23--4.25}; \HN{10.29--10.30}.
  \item Efectos de tratamiento con selección en observables (\emph{matching},
        IPW): ver D.1.
\end{itemize}

\subsection*{D.3 Econometría espacial}
Sigue el programa ECO224 (Estadística Aplicada: Análisis Espacial y
Econometría Espacial, 2020-2). Momento sugerido: después del Bloque 10, porque
los modelos espaciales se estiman por máxima verosimilitud.
\begin{itemize}[leftmargin=1.4em,itemsep=1pt]
  \item Matrices de pesos espaciales, I de Moran y LISA.
  \item Modelos SAR, SEM y SDM; efectos directos, indirectos y totales.
  \item Estimación por máxima verosimilitud.
  \item Lecturas: Anselin (2013); LeSage y Pace (2009); Elhorst (2014); Arbia (2014).
\end{itemize}

\subsection*{D.4 Econometría financiera}
Sigue los programas FIN214 (Econometría Financiera, 2025-2) y FIN222
(Tópicos de Econometría Financiera, 2023-2). La parte de series de tiempo
(ARIMA, VAR, cointegración, ARCH) requiere la carpeta Time Series.
\begin{itemize}[leftmargin=1.4em,itemsep=1pt]
  \item Elección discreta y modelos de clasificación para riesgo de crédito:
        Anderson (2007); Bloque 12.
  \item Aprendizaje automático en finanzas: Kelly y Xiu (2023); Athey e
        Imbens (2019).
  \item Econometría de los mercados financieros: Campbell, Lo y MacKinlay
        (1997), que está en la carpeta; Tsay (2010).
\end{itemize}

\subsection*{D.5 La regresión vista desde la estadística}
Sigue el programa EST225 (Análisis de Regresión, Facultad de Ciencias e
Ingeniería, 2025-2). Puede leerse en paralelo a los Bloques 2 y 3: añade el
diagnóstico de residuos, las medidas de influencia, la regresión polinómica,
la selección de variables, el análisis de varianza y el de covarianza.
Lecturas: Montgomery, Peck y Vining (2002); Fox (2008); Faraway (2014).

\subsection*{D.6 Proyecto empírico opcional}
Consolida los Bloques 8, 12 y 13 con datos peruanos. Momento sugerido:
después del Bloque 12; no tiene horas asignadas.
\begin{enumerate}[leftmargin=1.6em,itemsep=2pt]
  \item \textbf{Pregunta y datos.} Retornos de la educación o participación
        laboral femenina con la ENAHO (\WO{19.1--19.3}).
  \item \textbf{Diseño muestral.} Declarar estratos, conglomerados y factores
        de expansión (D.2).
  \item \textbf{Estimación.} MCO, VI (temas 8.2--8.4), probit (tema 12.1) y
        corrección de Heckman (tema 12.5).
  \item \textbf{Inferencia.} Errores agrupados y bootstrap (temas 5.4 y 7.4).
  \item \textbf{Redacción} del informe (\WO{19.4--19.5}).
\end{enumerate}
Software: Kleiber y Zeileis (2008) para R, que está en la carpeta, y las
secciones de programación de Hansen.

%------------------------------------------------------------------------------
\section*{Anexo E. Referencias}
\addcontentsline{toc}{section}{Anexo E. Referencias}

\subsection*{E.1 Textos del plan}
{\small
\begin{itemize}[leftmargin=1.4em,itemsep=1pt,label={}]
\item Davidson, R. y J. G. MacKinnon (2004). \emph{Econometric Theory and Methods}. Oxford University Press. Edición de los autores, 2021.
\item Greene, W. H. (2012). \emph{Econometric Analysis}, 7.ª ed. Pearson.
\item Hansen, B. E. (2021). \emph{Econometrics}. Manuscrito, revisión del 11 de marzo de 2021, University of Wisconsin. Publicado por Princeton University Press, 2022.
\item Hayashi, F. (2000). \emph{Econometrics}. Princeton University Press.
\item Wooldridge, J. M. (2012). \emph{Introductory Econometrics: A Modern Approach}, 5.ª ed. Cengage. Traducción: \emph{Introducción a la econometría. Un enfoque moderno}, 4.ª ed., Cengage, 2009.
\end{itemize}
}

\subsection*{E.2 Libros de consulta de la carpeta (sin horas)}
{\small
\begin{itemize}[leftmargin=1.4em,itemsep=1pt,label={}]
\item Angrist, J. D. y J.-S. Pischke (2009). \emph{Mostly Harmless Econometrics: An Empiricist's Companion}. Princeton University Press.
\item Cameron, A. C. y P. K. Trivedi (2005). \emph{Microeconometrics: Methods and Applications}. Cambridge University Press.
\item Campbell, J. Y., A. W. Lo y A. C. MacKinlay (1997). \emph{The Econometrics of Financial Markets}. Princeton University Press.
\item Davidson, R. y J. G. MacKinnon (1993). \emph{Estimation and Inference in Econometrics}. Oxford University Press.
\item Hansen, B. E. (2022). \emph{Probability and Statistics for Economists}. Princeton University Press.
\item Hastie, T., R. Tibshirani y J. Friedman (2009). \emph{The Elements of Statistical Learning}, 2.ª ed. Springer.
\item Huber, M. \emph{Causal Analysis: Impact Evaluation and Causal Machine Learning}. Universidad de Friburgo (versión de la carpeta).
\item Kleiber, C. y A. Zeileis (2008). \emph{Applied Econometrics with R}. Springer.
\item White, H. (2001). \emph{Asymptotic Theory for Econometricians}, ed. revisada. Academic Press.
\item Wooldridge, J. M. (2002). \emph{Econometric Analysis of Cross Section and Panel Data}. MIT Press.
\end{itemize}
}

\subsection*{E.3 Referencias tomadas de los programas de la PUCP}
{\small
\begin{itemize}[leftmargin=1.4em,itemsep=1pt,label={}]
\item Abadie, A., S. Athey, G. Imbens y J. Wooldridge (2017). When Should You Adjust Standard Errors for Clustering? NBER Working Paper 24003.
\item Amemiya, T. (1985). \emph{Advanced Econometrics}. Basil Blackwell.
\item Anderson, R. (2007). \emph{The Credit Scoring Toolkit}. Oxford University Press.
\item Anselin, L. (2013). \emph{Spatial Econometrics: Methods and Models}. Springer.
\item Arbia, G. (2014). \emph{A Primer for Spatial Econometrics: With Applications in R}. Springer.
\item Athey, S. y G. Imbens (2019). Machine Learning Methods That Economists Should Know About. \emph{Annual Review of Economics}, 11, 685--725.
\item Bajari, P. et al. (2021). Hedonic Prices and Quality Adjusted Price Indices Powered by AI. Cemmap.
\item Cameron, A. C. y D. L. Miller (2015). A Practitioner's Guide to Cluster-Robust Inference. \emph{Journal of Human Resources}, 50(2), 317--372.
\item Casella, G. y R. L. Berger (2021). \emph{Statistical Inference}. Cengage.
\item Chernozhukov, V., D. Chetverikov, M. Demirer, E. Duflo, C. Hansen, W. Newey y J. Robins (2018a). Double/Debiased Machine Learning for Treatment and Structural Parameters. \emph{The Econometrics Journal}, 21(1).
\item Chernozhukov, V., M. Demirer, E. Duflo e I. Fernández-Val (2018b). Generic Machine Learning Inference on Heterogeneous Treatment Effects in Randomized Experiments. NBER Working Paper 24678.
\item Chernozhukov, V., C. Hansen y M. Spindler (2015). Valid Post-Selection and Post-Regularization Inference: An Elementary, General Approach. \emph{Annual Review of Economics}, 7, 649--688.
\item Elhorst, J. P. (2014). \emph{Spatial Econometrics: From Cross-Sectional Data to Spatial Panels}. Springer.
\item Faraway, J. (2014). \emph{Linear Models with R}. Chapman \& Hall.
\item Fox, J. (2008). \emph{Applied Regression Analysis and Generalized Linear Models}. Sage.
\item García, L. (2022). \emph{Econometría 1}, 3.ª ed. Fondo Editorial PUCP.
\item Groves, R. M. et al. (2011). \emph{Survey Methodology}. Wiley.
\item Heckman, J. y R. Pinto (2015). Causal Analysis after Haavelmo. \emph{Econometric Theory}, 31(1), 115--151.
\item Imbens, G. W. y D. B. Rubin (2015). \emph{Causal Inference for Statistics, Social, and Biomedical Sciences}. Cambridge University Press.
\item Kelly, B. y D. Xiu (2023). Financial Machine Learning. NBER Working Paper 31502.
\item Lehmann, E. L. y J. P. Romano (2005). \emph{Testing Statistical Hypotheses}, 3.ª ed. Springer.
\item LeSage, J. y R. K. Pace (2009). \emph{Introduction to Spatial Econometrics}. CRC Press.
\item Lumley, T. (2010). \emph{Complex Surveys: A Guide to Analysis Using R}. Wiley.
\item Montgomery, D., E. Peck y G. Vining (2002). \emph{Introducción al análisis de regresión lineal}. Compañía Editorial Continental.
\item Solon, G., S. J. Haider y J. M. Wooldridge (2015). What Are We Weighting For? \emph{Journal of Human Resources}, 50(2), 301--316.
\item Tsay, R. S. (2010). \emph{Analysis of Financial Time Series}, 3.ª ed. Wiley.
\item van der Vaart, A. W. (2000). \emph{Asymptotic Statistics}. Cambridge University Press.
\end{itemize}
}

\vspace{1.2cm}
\franjas[2pt]
\begin{center}\small\color{tonoPizarra}
Plan de lectura de elaboración propia. Los títulos de los capítulos
corresponden a las obras citadas y se reproducen únicamente con fines de
referencia bibliográfica. La asignación de capítulos y secciones se basa en
los índices de los libros de la carpeta; los temas de la PUCP proceden de los
programas citados en los Anexos C y D.
\end{center}

\end{document}
